%% ======================================================================
%% GovSmart -- Kelompok 6 -- Workshop Pemrograman Framework
%% Format : Springer LNCS (llncs.cls + splncs04 reference style), as in format.pdf
%% Build  : Overleaf -> upload Springer's "llncs.cls" with this file, pdfLaTeX.
%%          Run pdflatex twice (references are inline, no BibTeX needed).
%%
%% HOW TO USE THIS FILE
%%   \res{...}  (blue)  = number/text that must come from the pilot data.
%%   \todo{...} (red)   = decision or fact the authors must supply.
%%   Before submission: search for "\res{" and "\todo{" -> none may remain.
%%   Never fill a \res{} with an estimate. Guide: "Conclusion ... angka hasil nyata".
%%
%% IF NO OBSERVATION EVIDENCE EXISTS when the Introduction is finalised,
%% replace paragraph 2 by the fallback wording given in the comment there.
%% ======================================================================
\documentclass[runningheads]{llncs}

\usepackage[T1]{fontenc}
\usepackage{graphicx}
\usepackage{booktabs}
\usepackage{array}
\usepackage{xcolor}
\usepackage{url}

% ---- drafting helpers (remove all uses before submission) ----
\newcommand{\res}[1]{\textcolor{blue}{[#1]}}
\newcommand{\todo}[1]{\textcolor{red}{\textbf{[TODO: #1]}}}
% ---- status tokens exactly as the portal displays them ----
\newcommand{\st}[1]{\textit{#1}}

\begin{document}

\title{Evaluating Complaint Submission and Status Comprehension in a Citizen Service Portal}
\titlerunning{Complaint Submission and Status Comprehension in GovSmart}

\author{Mochammad Zacki Ibrahim \and Mohammad Faiz Khoiru Izzuddin}
\authorrunning{M. Z. Ibrahim and M. F. K. Izzuddin}

\institute{Politeknik Elektronika Negeri Surabaya (PENS), PSDKU Lamongan,\\
D3 Informatics Engineering, Department of Informatics and Computer Engineering,\\
Lamongan, Indonesia\\
\email{\todo{author e-mail addresses}}}

\maketitle

%% ----------------------------------------------------------------------
%% Abstract: write LAST (problem - objective - method - result - conclusion).
%% Target 150-250 words once the \res{} fields are filled.
%% ----------------------------------------------------------------------
\begin{abstract}
A complaint portal helps only if citizens can send a complete report and
understand what the displayed status means for them. This paper evaluates both
abilities in GovSmart, a prototype portal for citizen complaints to a local
government. We ask whether users can submit a complete complaint and understand
the next action from the status shown. In a pilot with \res{N} prospective
users, each participant performed three tasks (submit, revise, and find the
status of a report) from a synthetic road-damage scenario and explained five
report statuses against a prepared answer key. We also tested input validation
and report ownership with two citizen accounts. Outcome, time, help, and errors
were logged per participant and task. \res{X of N submissions were complete
without help; Y of Z status interpretations were correct, and the most
misread term was T; V of W validation and ownership cases passed.} Failures
were traced to \res{specific form parts and status terms} and linked to
interface revisions. The pilot is small and uses synthetic reports, so it
supports problem identification only. It does not show higher citizen trust or
better service quality of any agency.
\keywords{Citizen complaint portal \and Usability evaluation \and Task success \and
Status comprehension \and E-government \and Form validation.}
\end{abstract}

%% ======================================================================
\section{Introduction}

%% --- Paragraph 1 (guide): citizen flow -- report, attach evidence, monitor ---
A citizen who notices a damaged public facility, such as a road, registers and
logs in to the portal. The citizen writes a title and a description, may add the
coordinates of the location, and may attach a photo as evidence. The portal
assigns the report to a department and an urgency level with keyword-based
routing and stores it with the status \st{diterima} (``received''). While the
report has this status its owner can edit or cancel it. Afterwards, department
staff update the status, may leave an officer note, and may upload a feedback
document that the citizen can open. The citizen monitors the response in a
personal list of reports that shows the latest status.

%% --- Paragraph 2 (guide): barriers observed in the first tasks, not assumed ---
Forms and status displays can create barriers for citizens, but we do not assume
that these barriers hold for all citizens. We ground the problem statement in
observations of the first tasks of our pilot: \res{describe the form fields at
which participants hesitated or erred and the status terms they misread, with
counts per field and per term}. Two design properties of the prototype motivated
the checks reported here. The status is shown to citizens as the lowercase
token returned by the API (for example \st{diterima}), and the evidence photo is
sent in a second request after the report itself has been saved, so a report can
exist without its attachment.
%% FALLBACK if no observation evidence exists (guide: state it as a prototype/scenario need,
%% not as an agency finding): "We treat the form and the status display of the prototype
%% as a scenario in which complete submission and correct status interpretation must be
%% verified; no field evidence about any agency's citizens is claimed."

%% --- Paragraph 3 (guide): portals, task success + status comprehension; vs SIGAP ---
Complaint portals have mostly been evaluated through satisfaction
surveys~\cite{yeh2010,undp2023}, qualitative implementation
studies~\cite{indrayani2025}, or field usability tests with a small number of
citizens~\cite{dahri2020}. Research on web forms studies completion and error
handling~\cite{wroblewski2008,inal2016}. In our search we found no study of a
complaint portal that asks participants to explain, in their own words, what a
status means and what they must do next, although keeping users informed about
system status is a long-standing usability heuristic~\cite{nielsen1994}. Our
work also differs from SIGAP, a companion prototype study in the same course that
examines the internal consistency of workflow transitions and assignment control
on the server side \todo{cite SIGAP as agreed with the lecturer}. We examine
instead what a citizen can do and understand with the portal.

%% --- Paragraph 4 (guide): contribution = citizen-side usability evaluation ---
This paper makes a citizen-side usability evaluation of the GovSmart portal. Its
contributions are: (i)~a task-based pilot on complete submission, revision, and
status finding; (ii)~a status-comprehension test against a written answer key;
and (iii)~a check of input validation and report ownership, with each proposed
interface improvement linked to an error that we observed. We answer one
research question: \emph{can users submit a complete complaint and understand the
next action from the status shown?}

%% ======================================================================
\section{Related Work}

\subsection{Evaluation of Citizen Complaint Portals}
The Kaohsiung Citizen Electronic Complaint System was assessed with satisfaction
and service quality as indicators, and the ability to resolve complaints
contributed most to overall satisfaction~\cite{yeh2010}. A user survey of the
Indonesian national channel SP4N-LAPOR! with 2,420 respondents reported, among
other indicators, that 90.8\% rated the simplification of complaint submission
positively~\cite{undp2023}. Both measure opinion after use rather than what users
can do. A qualitative study of the Lapak Aduan service in Banyumas, based on
observation, interviews, and documentation, found shortcomings in usefulness and
ease of use and low utilisation~\cite{indrayani2025}. Closest to our design,
field tests of the Pakistan Citizens' Portal app measured efficiency and
effectiveness following ISO 9241-11~\cite{iso9241} and satisfaction with a
questionnaire, with more than a dozen citizens, and located weaknesses in
registration and findability~\cite{dahri2020}.

\subsection{Form Completion}
Form guidelines ask that labels and input fields tell users what data is expected
and that errors are shown clearly~\cite{wroblewski2008}. An eye-tracking study of a
four-step web form with 32 participants compared four locations of error
messages; the researcher did not intervene during form filling, and the group
differences in completion time were not significant, while participants
preferred format restrictions, examples in error messages, and immediate
feedback~\cite{inal2016}. An earlier study with 303 participants found messages
placed next to the erroneous field most effective~\cite{seckler2012}. These
studies use generic forms, not complaint reports with evidence attachments.

\subsection{Status Comprehension}
Keeping users informed about what is going on is one of the usability
heuristics~\cite{nielsen1994}, but a status is only informative if its wording is
understood. Label-comprehension studies exist in other domains; for example,
studies of app privacy labels identified widely misunderstood technical
terms~\cite{keswani2025}. We did not find a comparable study for complaint
statuses, so we test it directly (Sect.~\ref{sec:status}).

\subsection{Comparison}
Table~\ref{tab:related} compares the studies on participants, tasks, moderator
help, and success metrics, the four aspects named in our evaluation plan.

\begin{table}[t]
\caption{Comparison of related studies. Entries follow the abstracts or summaries
of the sources; ``n/s'' = not stated there.}
\label{tab:related}
\centering
\footnotesize
\setlength{\tabcolsep}{3pt}
\begin{tabular}{@{}p{2.0cm}p{2.4cm}p{2.4cm}p{1.8cm}p{2.4cm}@{}}
\toprule
Study & Participants & Tasks & Moderator help & Success metrics \\
\midrule
Yeh and Chu~\cite{yeh2010} & Citizens (survey) & None (retrospective survey) & n/a & Satisfaction, service quality \\
UNDP survey~\cite{undp2023} & 2,420 respondents who used SP4N-LAPOR! & None (survey) & n/a & Share rating indicators positively \\
Indrayani et al.~\cite{indrayani2025} & n/s & Observation and interviews & n/s & Qualitative findings on usefulness, ease of use \\
Dahri et al.~\cite{dahri2020} & More than a dozen citizens & Field tests of app screens, incl.\ lodging a complaint & n/s & Efficiency, effectiveness (ISO 9241-11), satisfaction \\
Inal~\cite{inal2016} & 32 students & Four-step web form & None (no intervention) & Completion, error recognition, error correction time \\
\textbf{This study} & \res{N} prospective users & Submit, revise, find status; explain 5 statuses & Counted and coded per task & Success with/without help, time, errors, correct interpretations \\
\bottomrule
\end{tabular}
\end{table}

%% ======================================================================
\section{Methodology}

\subsection{System Under Evaluation}
GovSmart is a web portal with a React front end and a FastAPI back end on a MySQL
database. The front-end--back-end flow of submission, revision, and status display
was completed and frozen before the study (version \todo{commit hash and date}).
A report has a title of 5--100 characters, a description of 20--2,000
characters, an optional location (latitude in $[-90,90]$, longitude in
$[-180,180]$), and an optional photo (JPEG, PNG, or WEBP, at most 5~MB, sent in
a separate request). A new report receives the status \st{diterima}. A citizen
can edit or cancel a report only while it is \st{diterima}, can read only own
reports, and sees the officer note, when present, on the report card. The
personal report list returned by the API holds at most five reports.

\subsection{Participants, Consent, and Digital Experience}
We recruit \res{5--8} prospective users, \res{N} in the end, who are not members
of the project team. Each participant gives informed consent after hearing the
purpose, the recorded data, and the right to stop at any time. Participants are
identified as P1--P\res{N} only, and the study uses synthetic accounts and
synthetic reports, never real identity numbers. Before the tasks we record age
band, main device, how often the participant uses online government or other
transaction services, and any earlier experience of filing a complaint. We apply
the user-study protocol of the course guide \todo{add the protocol items from the
first pages of the guide, which were not in the file we used}.

\subsection{Scenario, Completeness Key, and Tasks}
\label{sec:tasks}
Every participant receives the same written scenario, without naming any button
or menu: a pothole has appeared on a road at a synthetic location, and a sample
photo is provided. The \emph{completeness key} for task~T1 lists what a complete
report must contain: a title that names the problem, a description that states
street, village or district, and the condition of the road, the sample photo
attached, and the correct department \todo{confirm the key items and whether the
coordinates are required}.

\paragraph{T1 -- Submit.} Report the problem to the local government.
\paragraph{T2 -- Revise.} After submitting, the participant is told that one
detail in the description is wrong and corrects it.
\paragraph{T3 -- Find status.} In an account that holds seeded reports in
different statuses, the participant finds the report named in the scenario and
states its current status.

A timer starts when the participant finishes reading the instruction and stops
when the correct result is reached, the time limit is reached, or the participant
gives up. Equivalent tasks use the same time limit (\todo{time limit per
task}). The moderator records every help given, using a fixed list of help types,
and every error. After each task we check the database for the data that were
actually stored (title, description, location, photo, department, status).
Outcomes are coded as success without help (S), success with help (SH), or
failure (F), where F includes a wrong or incomplete stored report.

\subsection{Status Comprehension Test}
\label{sec:status}
We choose the five statuses that exist in the portal and write the meaning and
the next action for each as the answer key (Table~\ref{tab:key}). The key is
agreed and frozen before the first session. Each status is shown on a concrete
report card, including the officer note where the status has one. The participant
explains in own words what the status means and whether any action is needed;
no answer options are offered, because options would leak the answer. The
participant then rates confidence from 1 to 5. Two team members code each answer
as correct, partly correct, or incorrect against the key and record the terms that
caused misreadings; differences are settled by discussion.

\begin{table}[t]
\caption{Draft answer key for the five statuses, derived from the behaviour of the
prototype. The wording is a proposal to be agreed and frozen by the authors.}
\label{tab:key}
\centering
\footnotesize
\setlength{\tabcolsep}{3pt}
\begin{tabular}{@{}p{1.5cm}p{4.0cm}p{2.2cm}p{3.4cm}@{}}
\toprule
Status shown & Intended meaning & Who is acting & Citizen's next action \\
\midrule
\st{diterima} & The system has received the report and routed it to a department; no officer has started work. & System, department queue & None required; the report can still be edited or cancelled. \\
\st{diproses} & A department officer is handling the report. & Department officer & Wait for an update; editing and cancelling are no longer possible. \\
\st{selesai} & The department has marked the report as handled. & No one & Read the officer note and open the feedback document, if any. \\
\st{ditolak} & The department declined the report; the reason is in the officer note. & No one & Read the reason; send a corrected new report if the problem remains. \\
\st{dibatalkan} & The citizen cancelled the report. & No one & None. \\
\bottomrule
\end{tabular}
\end{table}

\subsection{Validation and Ownership Test}
\label{sec:validation}
With citizen accounts A and B we submit the following cases and inspect the
message shown, the HTTP response, and the database: empty title and description;
title and description below and above their length limits; a file of a wrong type;
a file above 5~MB; an empty file; and the same submission sent twice, to check the
duplicate-submission policy \todo{define the intended policy before the test; a
report sent twice is either merged, rejected, or accepted as a second report}.
From account A we then open, edit, cancel, and attach a photo to the report of
account B through the interface and through the API. Expected result: invalid
input is rejected with a clear message, resubmission follows the policy, and the
report of B is not accessible to A.

\subsection{Data Recorded and Analysis}
Table~\ref{tab:measures} lists the measures. With fewer than ten participants we
report counts, medians, and the individual cases, and no significance tests. A
sample of this size is suited to finding usability problems rather than to
estimating rates~\cite{nielsen1993}. If we build an improved version, participants
perform equivalent but not identical tasks, and we discuss the learning effect;
an old-versus-new comparison is not treated as a strong causal experiment.

\begin{table}[t]
\caption{Measures recorded.}
\label{tab:measures}
\centering
\footnotesize
\begin{tabular}{@{}p{3.2cm}p{8.4cm}@{}}
\toprule
Measure & Definition \\
\midrule
Task outcome & S, SH, or F per participant and task (Sect.~\ref{sec:tasks}) \\
Time & Seconds from end of instruction to stop rule \\
Help and errors & Count and type per task, recorded by the moderator \\
Stored data & Fields actually stored versus the completeness key \\
Status interpretation & Correct, partly correct, or incorrect per status; misread terms \\
Confidence match & Confidence rating (1--5) compared with correctness \\
Validation and ownership & Pass or fail per case; duplicate reports; unauthorised changes; confusing messages \\
\bottomrule
\end{tabular}
\end{table}

%% ======================================================================
\section{Results and Discussion}
%% All numbers below come from the pilot logs. Leave nothing estimated.

\subsection{Data, Cases, and Conditions}
Table~\ref{tab:data} describes the participants and the experimental conditions.

\begin{table}[t]
\caption{Participants and conditions (output 1).}
\label{tab:data}
\centering
\footnotesize
\setlength{\tabcolsep}{3pt}
\begin{tabular}{@{}p{1.0cm}p{1.5cm}p{1.7cm}p{3.0cm}p{3.3cm}@{}}
\toprule
ID & Age band & Device & Online-service experience & Earlier complaint experience \\
\midrule
P1 & \res{} & \res{} & \res{} & \res{} \\
P2 & \res{} & \res{} & \res{} & \res{} \\
\multicolumn{5}{@{}p{11.2cm}@{}}{\res{one row per participant; add the session date, place, and time limit per task below the table}} \\
\bottomrule
\end{tabular}
\end{table}

\subsection{Results per Category}
Table~\ref{tab:results} reports the results per task, per status, and per
validation category.

\begin{table}[t]
\caption{Results per category (output 2). For tasks: S = success without help, SH = success with help, F = failure. For statuses: correct, partly correct, incorrect. For validation and ownership cases: pass or fail.}
\label{tab:results}
\centering
\footnotesize
\setlength{\tabcolsep}{4pt}
\begin{tabular}{@{}p{3.4cm}rrrrr@{}}
\toprule
Category & n & S/corr. & SH/part. & F/inc. & Med.\ time (s) \\
\midrule
T1 Submit & \res{} & \res{} & \res{} & \res{} & \res{} \\
T2 Revise & \res{} & \res{} & \res{} & \res{} & \res{} \\
T3 Find status & \res{} & \res{} & \res{} & \res{} & \res{} \\
\st{diterima} & \res{} & \res{} & \res{} & \res{} & -- \\
\st{diproses} & \res{} & \res{} & \res{} & \res{} & -- \\
\st{selesai} & \res{} & \res{} & \res{} & \res{} & -- \\
\st{ditolak} & \res{} & \res{} & \res{} & \res{} & -- \\
\st{dibatalkan} & \res{} & \res{} & \res{} & \res{} & -- \\
Invalid input cases & \res{} & \res{} & -- & \res{} & -- \\
Resubmission & \res{} & \res{} & -- & \res{} & -- \\
Cross-account access & \res{} & \res{} & -- & \res{} & -- \\
\bottomrule
\end{tabular}
\end{table}

\subsection{Failure Cases and Causes}
Table~\ref{tab:fail} lists the failures, each with its cause and the revision it led to.

\begin{table}[t]
\caption{Failure examples and causes (output 3).}
\label{tab:fail}
\centering
\footnotesize
\setlength{\tabcolsep}{3pt}
\begin{tabular}{@{}p{0.5cm}p{3.2cm}p{2.5cm}p{2.5cm}p{2.1cm}@{}}
\toprule
\# & Observation (anonymised) & Form part or status term & Cause & Revision \\
\midrule
1 & \res{} & \res{} & \res{} & \res{} \\
2 & \res{} & \res{} & \res{} & \res{} \\
\bottomrule
\end{tabular}
\end{table}

\subsection{Discussion}
%% Link every failure to a form part or a status term. Say why it happened.
\res{Answer the research question first: how many participants submitted a
complete report and found the right report without help, and how many understood
the next action for each status.} \res{Explain each failure through the form part
or status term in Table~\ref{tab:fail}, for example a label, a field limit, or a
lowercase status token, and say why participants read it as they did.}
The validation and ownership results are reported separately from the usability
results, because they concern the system rather than the participants:
\res{state which cases failed and what the database showed}. If an improved
version was tested, participants performed equivalent tasks, and part of any
improvement may come from learning; we therefore treat the comparison as
indicative only.

%% ======================================================================
\section{Conclusion}
\res{Answer the research question with the real numbers: X of N participants
submitted a complete report without help; Y of Z status interpretations were
correct, with T the most misread term; V of W validation and ownership cases
passed.} The pilot is small, uses synthetic reports and a convenience sample, and
was run on one prototype version. It therefore supports the identification of
problems in the form and in the status wording and the revisions linked to them.
It does not support conclusions about citizen trust or about the service quality
of any agency, because neither was measured.

\begin{credits}
\subsubsection{\ackname} The authors thank the pilot participants and the lecturer
of the course Workshop Pemrograman Framework \todo{confirm wording and name}.

\subsubsection{\discintname} The authors have no competing interests to declare that are
relevant to the content of this article.
\end{credits}

%% ======================================================================
%% References: numbered in order of first citation, splncs04 style.
%% Entries marked (verified) were checked against the publisher or journal page.
%% Entries marked (check) come from the authors' memory: confirm before submission.
\begin{thebibliography}{10}

\bibitem{yeh2010} % (verified)
Yeh, S.-C., Chu, P.-Y.: Evaluation of e-government services: a citizen-centric
approach to citizen e-complaint services. In: Citizens and e-Government:
Evaluating Policy and Management, pp. 400--417. IGI Global, Hershey (2010).
\url{https://doi.org/10.4018/978-1-61520-931-6.ch022}

\bibitem{undp2023} % (verified page; confirm publication date)
UNDP Indonesia: Survey reveals 8 out of 10 Indonesians recommend using SP4N-LAPOR!
system to report public service complaints. \url{https://www.undp.org/node/385696},
last accessed 2026/10/25

\bibitem{indrayani2025} % (verified)
Indrayani, E., Akbar, F., Diantoro, S.: Overcoming challenges in digital public
service delivery: insights from Lapak Aduan implementation. TRANSFORMASI: Jurnal
Manajemen Pemerintahan \textbf{17}(2), 103--121 (2025).
\url{https://doi.org/10.33701/jtp.v17i2.3677}

\bibitem{dahri2020} % (verified abstract; add page numbers from the PDF)
Dahri, A.S., Massan, S.-R., Maitlo, A.A.: Usability of eGovernance application for
citizens of Pakistan. 3C Tecnolog\'{\i}a, Special Issue, April 2020.
\url{https://3ciencias.com/articulos/articulo/usability-of-egovernance-application-for-citizens-of-pakistan/}

\bibitem{wroblewski2008} % (verified)
Wroblewski, L.: Web Form Design: Filling in the Blanks. Rosenfeld Media, New York
(2008)

\bibitem{inal2016} % (verified)
Inal, Y.: User-friendly locations of error messages in web forms: an eye tracking
study. Journal of Eye Movement Research \textbf{9}(5), 1--11 (2016).
\url{https://doi.org/10.16910/jemr.9.5.1}

\bibitem{nielsen1994} % (check)
Nielsen, J.: Enhancing the explanatory power of usability heuristics. In: Proceedings
of the SIGCHI Conference on Human Factors in Computing Systems (CHI '94),
pp. 152--158. ACM, New York (1994)

\bibitem{iso9241} % (check)
ISO 9241-11:2018: Ergonomics of human-system interaction -- Part 11: Usability:
Definitions and concepts. International Organization for Standardization, Geneva
(2018)

\bibitem{seckler2012} % (verified via the reference list of Inal 2016)
Seckler, M., Tuch, A.N., Opwis, K., Bargas-Avila, J.A.: User-friendly locations of
error messages in web forms: put them on the right side of the erroneous input
field. Interacting with Computers \textbf{24}(3), 107--118 (2012)

\bibitem{keswani2025} % (verified abstract; add pages)
Keswani, I., Walker, K., Clement, A., Kitur, E., Wonghirundacha, N., Aubrey, R.,
Song, V., Birrell, E.: User understandings of technical terms in app privacy
labels. In: USENIX Symposium on Usable Privacy and Security (SOUPS 2025). USENIX
Association (2025)

\bibitem{nielsen1993} % (check)
Nielsen, J., Landauer, T.K.: A mathematical model of the finding of usability
problems. In: Proceedings of the INTERACT '93 and CHI '93 Conference on Human
Factors in Computing Systems, pp. 206--213. ACM, New York (1993)

\end{thebibliography}

\end{document}